% !TEX root = ../main.tex
\chapter{Credit Default Swaps}\label{ch:cds}
\begin{paracol}{2}

\begin{sources}
\begin{tabularx}{\linewidth}{@{}l l L@{}}
\toprule
\textbf{Segment} & \textbf{Length} & \textbf{Content}\\
\midrule
\seg{20}{1}{L20.1 FT: Internationalization of the Euro} & 5:07 & Offshore euro trading in London; bonds replacing syndicated loans.\\
\seg{20}{2}{L20.2 Credit indices} & 3:40 & The Markit CDX high-yield index.\\
\seg{20}{3}{L20.3 Fischer Black (1970), risk-free security} & 2:14 & A risky bond $=$ risk-free bond $+$ interest rate risk $+$ default risk.\\
\seg{20}{4}{L20.4 What is a Credit Default Swap (CDS)} & 5:12 & Buyer and seller of protection on the balance sheet.\\
\seg{20}{5}{L20.5 Corporate bonds} & 3:43 & Price as present value; spreads and ratings.\\
\seg{20}{6}{L20.6 CDS pricing} & 11:24 & Bond $+$ IRS $=$ LIBOR $+s$; CDS premium $u$; arbitrage $s=u$; is it insurance?\\
\seg{20}{7}{L20.7 Market making in CDS} & 4:03 & Dealers hedge single names with the index.\\
\seg{20}{8}{L20.8 Example: Negative basis trade and Liquidity Risk} & 10:20 & UBS: AAA CDO tranches, insured, funded in the money market.\\
\seg{20}{9}{L20.9 Example: Private backstop of marketmaking in CDS} & 15:15 & AIG, Goldman Sachs and collateral calls: ``liquidity kills you quick''.\\
\seg{20}{10}{L20.10 Example: Synthetic CDO as Collateral Prepayment} & 6:43 & Paulson, Abacus and IKB.\\
\bottomrule
\end{tabularx}

\smallskip
\textbf{Files.} Transcripts \file{materials/transcripts/L20-P*.txt}; L20.1 has no YouTube captions and was transcribed locally from the audio (faster-whisper). Mehrling's notes \mn{20}{1}.
\end{sources}

\section*{Motivation}
All semester Mehrling promised that the course's balance sheets would handle credit default swaps; now he delivers \ts{20}{2}{0:06}. Think of a corporate bond with a lot
of credit risk: a credit default swap peels that risk off and sells it separately, just as an interest rate swap peels off interest rate exposure
\ts{20}{2}{3:11}. Three examples from the crisis (UBS, AIG, Abacus) then show that the trouble came not from defaults but from liquidity.

% ------------------------------------------------------------------
\section{FT: the internationalization of the euro}\label{sec:cds-ft}

Mehrling had just come back from Rome, where Italian central bankers and securities regulators asked his shadow banking team whether shadow banking could channel
credit to small and medium-sized enterprises: they depend on banks, and banks, having bought government bonds, now need to build up capital and are not lending
\ts{20}{1}{0:00}. He reads two articles of the day as symptoms of sea changes in the European financial system
\ts{20}{1}{1:24}.

\begin{casestudy}[FT, November 2012: two European evolutions]
\begin{enumerate}
  \item ``France targets UK euro trade supremacy'': 40\% of euro trading takes place in London, and France wants it inside the eurozone, under control
        \ts{20}{1}{1:35}. Mehrling recalls the rise of the Eurodollar market in London and the old American worry about losing control of the dollar;
        the US concluded it could not stop it, and the same holds for the euro: people who want to trade euro derivatives will find somewhere to do it, London or
        Hong Kong. A ``little naive'', but a sign of the maturation of the euro as an international currency \ts{20}{1}{2:03};
  \item ``Syndicated bank loans in eurozone slump to ten-year low'': top-rated corporations no longer take syndicated loans (a bank making a loan much bigger than it
        could hold and parcelling out pieces to other banks) but issue bonds directly, in big amounts: the capital market replaces bank lending, as in the US
        evolution from indirect to direct finance (lecture~\ref{ch:directindirect}) \ts{20}{1}{3:21}.
\end{enumerate}
One evolution is in the money market (London), the other in the capital market; link the two and you have shadow banking, money market funding of capital market
lending, driven by market forces, while ``none of the regulations are right for this'' \ts{20}{1}{4:15}.
\end{casestudy}

\begin{coursenote}
The figure of 40\% is as said in class, quoting the FT. The ``30 years ago'' of the Eurodollar debate is approximate: the Eurodollar market grew in London from the
late 1950s, and the US debate on controlling it was intense in the late 1960s and 1970s \ts{20}{1}{1:46}.
\end{coursenote}

% ------------------------------------------------------------------
\section{Credit indices}\label{sec:cds-index}

\begin{finance}[The Markit CDX family]
Markit's CDX indices are, in its words, the standard tradable family of credit default swap indices for North America and emerging markets: they offer exposure to
a diversified portfolio of credits, a hedge of credit risk, and liquidity that flows into the single-name CDS market \ts{20}{2}{0:27}.
The CDX.NA.HY (North America high yield, i.e.\ junk) index: spreads of 450--580 basis points over a riskless rate between late September and the day of the lecture;
the index is set at 100 and moves inversely to the spread (about 530 bp at that point). Its constituents are 100 companies, each with 1\% weight, with ratings
such as BB, B, CCC \ts{20}{2}{1:29}.
\end{finance}

% ------------------------------------------------------------------
\section{Fischer Black's vision (1970)}\label{sec:cds-black}

In an unpublished 1970 memo, which Mehrling read while writing Black's biography, the young Fischer Black imagined a long-term corporate bond sold to three
separate persons: one supplying the money, one bearing the interest rate risk and one the default risk, the last two putting up no capital but perhaps posting
collateral. ``That's the world that came to be'' \ts{20}{3}{0:05}.

\begin{keyformulas}
\textbf{Black's decomposition} \ts{20}{3}{0:48}:
\[
  \text{risk-free security} = \text{risky (corporate) bond} + \text{interest rate hedge (IRS)} + \text{credit hedge (CDS)}.
\]
\end{keyformulas}

\begin{history}[Fischer Black]
Fischer Black (1938--1995), co-author of the Black--Scholes option pricing formula (1973), worked at Arthur D.~Little, the University of Chicago, MIT and Goldman
Sachs. Mehrling's biography draws on his unpublished papers, including the 1970 memo ``Fundamentals of Liquidity''.\par
\emph{Reference:} P.~Mehrling, \emph{Fischer Black and the Revolutionary Idea of Finance}, Wiley, 2005.
\end{history}

% ------------------------------------------------------------------
\section{What is a credit default swap}\label{sec:cds-what}

There is a buyer and a seller of ``insurance''. The buyer is \emph{long the swap} and \emph{short credit risk} \ts{20}{4}{0:04}. Balance sheet
construction \ts{20}{4}{0:56}:
\begin{enumerate}
  \item start from a risky corporate bond;
  \item add (in brackets: notional) a long Treasury bond and a short position in the same corporate bond, of the same maturity: the net exposure is now a
        Treasury bond. The pair is the CDS: it carves off the credit risk, the 530 bp;
  \item add a short Treasury bond and a long Treasury bill: the net exposure is a Treasury bill. This pair, long fixed and short floating, is an interest rate swap.
\end{enumerate}

\begin{nfigure}
\begin{taccts}
\tacct[2.6cm]{Buyer of protection}{Corporate bond\\{}[$+$Treasury bond, $-$corporate bond] $=$ CDS (long swap)\\{}[$+$T-bill, $-$Treasury bond] $=$ IRS}{}\qquad
\tacct[2.6cm]{Seller of protection}{}{CDS (short swap, long credit risk)\\IRS}
\end{taccts}
\caption{A CDS and an IRS as notional swaps of IOUs: the hedges are booked as contingent assets of the buyer and contingent liabilities of the seller
\ts{20}{4}{2:50}.}\label{fig:cds-bs}
\end{nfigure}

\begin{definition}[New concepts: credit default swap, protection buyer and seller]
A \emph{credit default swap} (CDS) is a contract in which the \emph{protection buyer} pays a periodic premium and, if the reference bond defaults, receives
from the \emph{protection seller} the difference between face value and the bond's post-default value \ts{20}{4}{0:04}\ts{20}{6}{4:28}.
(\emph{Default}: from Old French \emph{defaute}, a failing or lack.)
\end{definition}

The course books any hedge as a contingent asset, so the buyer lists both the long credit exposure (the bond) and the short one (the CDS) as assets; real-world
practice is not consistent, often off balance sheet \ts{20}{4}{3:12}. The seller is long credit risk and short the swap: the
lingo is consistent between CDS and IRS. ``Credit default swaps are not some weird exotic creature of Wall Street'' \ts{20}{4}{4:33}.

% ------------------------------------------------------------------
\section{Corporate bonds}\label{sec:cds-bonds}

For the first time in the course, a corporate bond: a fixed coupon every six months or year for, say, ten years, and the principal at the end; unlike a mortgage it
is not amortized \ts{20}{5}{0:04}.
\begin{keyformulas}
\textbf{Bond price} \ts{20}{5}{1:05}: with coupons $C_t$, face value $F$ at maturity $T$, and a risk-weighted discount rate
(the risk-free rate plus a credit spread $\delta$, e.g.\ 530 bp),
\[
  P_0 = \sum_{t=1}^{T}\frac{C_t}{(1+r+\delta)^t} + \frac{F}{(1+r+\delta)^T}.
\]
\end{keyformulas}
Bonds have ratings, and at any moment there are credit spreads between rating classes, which change over time. So a bond's price changes when the spread for its
rating changes and when its rating changes, a ``double whammy'' if a downgrade comes with wider spreads \ts{20}{5}{2:12}.

% ------------------------------------------------------------------
\section{CDS pricing}\label{sec:cds-pricing}

\paragraph{Step 1: strip the interest rate risk.} Take a ten-year bond with constant coupon and swap fixed for LIBOR: what remains is a ten-year floating-rate bond paying
LIBOR $+\,s$, with the credit spread $s$ locked in \ts{20}{6}{0:02}.

\paragraph{Step 2: the CDS as a swap.} Each year the bond survives, one side pays LIBOR times face value and the other LIBOR $+\,u$: the protection buyer pays the net
premium $u$ times face value \ts{20}{6}{1:51}. If the bond defaults, say in period five, the buyer
swaps it for a Treasury bond: it receives face value minus the bond's price at that time \ts{20}{6}{4:05}. The contract follows the
two parts of a bond: the annual (coupon) business and the terminal payment if it defaults \ts{20}{6}{5:35}.

\begin{nfigure}
\begin{tikzpicture}[font=\scriptsize,x=1.3cm,y=1cm]
  \draw[-{Stealth[length=4pt]}] (0,0) -- (6.4,0) node[right]{period};
  \foreach \t in {1,...,4} {\draw[thick,-{Stealth[length=3pt]}] (\t,0) -- (\t,-0.7); \node[below] at (\t,-0.7) {$-u F$};}
  \draw[very thick,-{Stealth[length=4pt]}] (5,0) -- (5,1.8);
  \node[above,align=center] at (5,1.8) {default: $+\,(F - P_5)$};
  \foreach \t in {1,...,5} \node[above] at (\t,0.05) {\t};
\end{tikzpicture}
\caption{Cash flows to the protection buyer: the premium $u$ times face value $F$ while the bond survives, face value minus the bond's price $P_5$ at default
\ts{20}{6}{3:45}.}\label{fig:cds-flows}
\end{nfigure}

\begin{keyformulas}
\textbf{CDS pricing by arbitrage} \ts{20}{6}{6:36}: the CDS premium $u$ should equal the bond's credit spread $s$;
otherwise go long the bond and insure it (if $u<s$), or short the bond and sell insurance (if $u>s$). This is how CDS are priced, ``assuming that there's no
liquidity premium anywhere in this picture''. The difference $u-s$ is the \emph{CDS--bond basis}.
\end{keyformulas}

The notes call the protection seller's IOU a \emph{mirror bond}: it promises LIBOR $+U$ on face value as long as the issuer pays, and the liquidation value on default; the buyer's IOU promises
LIBOR, and face value on default. On default the buyer recovers $F$ in full: $P_5$ from the bond, $F-P_5$ from the swap. $U$ is the premium that makes the two IOUs equal in value at inception; but
since any change in the spread $S$ moves its value, ``CDS is not so much insurance against eventual default as it is insurance against change in the credit spread'' \mn{20}{3}\mn{20}{4}. Pooling
CDS reduces idiosyncratic risk (the less correlated, the better), letting sellers charge less: a flourishing swap market should lower the price of credit risk overall, ``and that is exactly what
happened''; for subprime the index used for hedging was ABX, the index of top-tranche mortgage CDOs \mn{20}{1}\mn{20}{4}\mn{20}{5}.

\paragraph{Not just insurance.} A CDS is a promise to pay a fixed premium for its life, but its value fluctuates, inversely to the bond, whenever the warranted
discount rate changes; so CDS are used not only as insurance against default but as a way to get exposure to credit risk, traded and reversed like any liquid
derivative. Unlike fire insurance, you are not just waiting for the house to burn \ts{20}{6}{8:00}.

\begin{intuition}[Is a CDS insurance?]
``Insurance'' is used to remember which side of the balance sheet to book; but many people dislike the word, on both sides of the political spectrum. If it were
insurance it would be regulated as insurance, with reserves and suitability standards; and insurance works by pooling independent risks (fires), while credit risk
is systemic: corporate bonds and their spreads move together. The mathematicians and theorists say: it is not insurance, just a swap
\ts{20}{6}{9:43}.
\end{intuition}

% ------------------------------------------------------------------
\section{Market making in CDS}\label{sec:cds-dealers}

Opposite the protection buyer is typically a dealer writing CDS on bonds $i$, $j$, $k$\ldots: a portfolio of CDS is like a diversified portfolio of corporate bonds
\ts{20}{7}{0:03}. The dealer cannot find someone wanting exactly the opposite single-name positions, so it hedges with a CDS on the
CDX index, which is correlated with its portfolio; on the other side may be an asset manager seeking exposure to the bond market as a whole
\ts{20}{7}{1:08}.

\begin{nfigure}
\begin{taccts}
\tacct[1.9cm]{Hedger (single name)}{CDS on bond $i$ (long)}{}\quad
\tacct[1.9cm]{CDS dealer (investment bank)}{CDS on CDX index}{CDS on bonds $i, j, k, \ldots$}\quad
\tacct[1.9cm]{Asset manager}{}{CDS on CDX index (long credit risk)}
\end{taccts}
\caption{A CDS dealer hedging a portfolio of single names with the index \ts{20}{7}{1:31}.}\label{fig:cds-dealer}
\end{nfigure}

With a dealer in the middle, order flow pushes CDS prices around without necessarily moving bond prices; someone must arbitrage to bring them in line. So
arbitrage opportunities appear whenever CDS prices are out of line with bond prices, ``and sometimes it is'' one \ts{20}{7}{2:55}.

% ------------------------------------------------------------------
\section{UBS: the negative basis trade and liquidity risk}\label{sec:cds-ubs}

UBS lost a lot of money on what it thought was a riskless trade, and its regulator made it write a public report explaining how; Mehrling wishes US regulators had
required the same of Citibank and the TARP banks \ts{20}{8}{0:03}. UBS noticed that the spread $s$ on a security exceeded the
cost $u$ of insuring it against default, at least from AIG \ts{20}{8}{0:45}:
\begin{enumerate}
  \item it bought AAA tranches of CDOs on mortgage-backed securities, some on subprime mortgages, with the credit risk supposedly carved off into lower tranches
        \ts{20}{8}{1:26};
  \item its European regulators required insurance to use them as bank collateral, so it bought CDS, from AIG among others: ``double insurance''
        \ts{20}{8}{2:09};
  \item it funded them in the money market, asset-backed commercial paper or repo, with the tranche as collateral: money market funding of capital market lending,
        shadow banking, but right on UBS's balance sheet, unlike Citibank's structured investment vehicles off balance sheet \ts{20}{8}{3:17}.
\end{enumerate}

\begin{nfigure}
\begin{taccts}
\tacct[3.0cm]{UBS (negative basis trade)}{AAA CDO tranche (yield LIBOR $+\,s$)\\CDS bought from AIG (premium $u<s$)}{Money market funding (ABCP, repo), collateral: the tranche}
\end{taccts}
\caption{The negative basis trade \ts{20}{8}{1:48}.}\label{fig:cds-ubs}
\end{nfigure}

\paragraph{The numbers.} The notes give them from UBS's report: UBS paid AIG 11 basis points for 100\% protection on a super-senior CDO tranche, financed the tranche in the wholesale money
market, and netted an apparently riskless 20 basis points; ``because it was apparently riskfree, they did massive amounts of it'' \mn{20}{5}. In the notes' notation, if the CDS premium $U$ is below
the spread $S$, the buyer of protection swaps out the credit risk and keeps $S-U$ over LIBOR.

\paragraph{Booking 20 years of profit today.} The \emph{negative basis trade} left money over each period for the life of the tranche, 10 or 20 years; UBS booked the
whole stream as profit today and paid bonuses on it \ts{20}{8}{4:42}.

\paragraph{The liquidity spiral.} Then small doubts about the AAA tranche: its value fluctuated a little; the CDS moved the other way, fine; but the tranche, not the
CDS, was the collateral, so the funding could not be rolled, the tranche had to be sold, pushing its price down further, and others doing the same trade made a
downward liquidity spiral \ts{20}{8}{5:47}. Moreover, AIG's insurance was expensive, so late in the game UBS insured
only against a 2\% fall (from 100 to 98) and kept the tail risk itself \ts{20}{8}{6:48}. They felt they took no risk at all and lost
billions, because they assumed they could always roll their funding \ts{20}{8}{7:50}.

\begin{coursenote}
A student asks why the lenders did not accept the CDS as collateral together with the tranche. Perhaps by analogy with Bagehot's bills, where the acceptance was
attached to the bill; here they are separate securities, and money market funds are restricted in the collateral they may accept. Mehrling recalls that when AIG
was failing he proposed that the government write CDS itself, since a CDS is like an acceptance, which is what a liquidity problem needs
\ts{20}{8}{8:11}.
\end{coursenote}

\begin{history}[The UBS report]
At the request of the Swiss Federal Banking Commission, UBS published in April 2008 a ``Shareholder Report on UBS's Write-Downs'', describing among other things its
negative basis trades and the ``Amplified Mortgage Portfolio'' (AMPS), in which only a small first slice of losses on super-senior positions was hedged.\par
\emph{Reference:} UBS AG, \emph{Shareholder Report on UBS's Write-Downs}, 18 April 2008.
\end{history}

% ------------------------------------------------------------------
\section{AIG and Goldman Sachs: a private backstop}\label{sec:cds-aig}

Goldman Sachs, a dealer in CDS selling protection to its clients (not completely matched book: short credit risk when it mattered), kept buying CDS from AIG, which held
CDS liabilities and capital \ts{20}{9}{0:00}. When the CDS came into the money, AIG's contingent liability, zero when
written, became, say, \$30 billion owed to Goldman \ts{20}{9}{1:36}. Goldman had written the contracts so that if AIG's rating fell below AAA it
would post margin; Moody's downgraded AIG, and it had to \emph{pay} the \$30 billion, not just owe it \ts{20}{9}{2:19}.

\begin{nfigure}
\begin{taccts}
\tacct[2.4cm]{Goldman Sachs (dealer)}{CDS bought from AIG\\($+$\$30 bn when in the money: margin received)}{CDS sold to clients\\($+$\$30 bn owed to clients)}\qquad
\tacct[2.4cm]{AIG}{Capital (bonds)\\($-$cash paid as margin)}{CDS on AAA securities\\($+$\$30 bn contingent liability)}
\end{taccts}
\caption{AIG as private backstop of CDS market making; ``say \$30 billion'' in class; the notes give about \$30 billion of collateral posted to Goldman's segregated account by the
time AIG failed \ts{20}{9}{1:16}\mn{20}{6}.}\label{fig:cds-aig}
\end{nfigure}

\begin{intuition}[Liquidity kills you quick]
AIG thought it was insuring against defaults that would never happen; but you can go bankrupt without any default if the CDS fluctuate in value and you must pay
margin. The contractual promise to post collateral killed AIG; monoline insurers that wrote CDS without margin saw the marks go against them but paid nothing.
``Insolvency you can survive for a long time\ldots{} but liquidity kills you quick'': AIG, the anchor of the system, fell two or three days after Lehman in September
2008, and the system went into free fall \ts{20}{9}{2:40}.
\end{intuition}

\begin{coursenote}
The AIG executive is not named in class. Joseph Cassano, head of AIG Financial Products, told investors in August 2007 that it was hard to see a scenario in which
the unit would lose ``one dollar'' on these transactions (``a nickel'' in class). The US government took a 79.9\% equity stake in AIG in September 2008 (``80'' in
class).\par
\emph{References:} Financial Crisis Inquiry Commission, \emph{Final Report}, 2011, ch.~13; SIGTARP, \emph{Factors Affecting Efforts to Limit Payments to AIG Counterparties},
November 2009.
\end{coursenote}

\begin{intuition}[Was Goldman paid at par by the government?]
The notes correct ``a lot of loose talk'' \mn{20}{6}: since the CDS were marked to market, Goldman already held the collateral, it had already been paid. When AIG could not meet further collateral
calls the government took over, lending \$85 billion; that money was used to acquire the referenced securities at liquidation value to end the swaps. In the algebra of the pricing example, AIG had
already paid $F-P_5$; the government lent AIG $P_5$ to buy the bonds, which went onto the Fed's balance sheet as Maiden Lane II and III.
\end{intuition}

\paragraph{Should the government have clawed it back?} Neil Barofsky (SIGTARP), whose book is on the final-exam reading, wrote a report on AIG: Goldman and, even bigger,
Soci\'et\'e G\'en\'erale sucked the capital out of AIG, the government took over AIG, and the question was what to do about the money already paid to Goldman
\ts{20}{9}{4:03}. But to the extent Goldman was a market maker with offsetting positions, taking back the \$30 billion would have left it
unable to pay its clients and the whole thing unravels; ``I don't say that Geithner necessarily did the right thing'', but it was not simply unearned profit: the
point of a CDS is that when the corporate bond goes south you replace it with a Treasury \ts{20}{9}{5:05}.

\paragraph{No fraud, a misunderstanding of risk.} Mehrling is not sure there was fraud, at AIG or at UBS: a systematic misunderstanding of risk, a supposedly risk-free
arbitrage on which you bet the bank, and do more of it. ``If it looks like free money, look again'', with the money view: is the arbitrage just a liquidity premium
you think should be zero? Then you are becoming the lender of last resort to the whole system without being a central bank that can print money
\ts{20}{9}{7:28}.

\paragraph{Moral hazard?} A student suggests they took the risk expecting a bailout. There is always some such calculation, but Mehrling thinks they were surprised.
The one really relying on the government was Goldman: it chose AIG though it was more expensive (UBS found it too expensive), because AIG promised margin and had a
huge capital base; and AIG, reserving nothing and paying bonuses on each sale, could not be allowed to fail, with fire insurance policies worldwide. But a
counterparty that dies is not the right counterparty: Goldman surely did not expect to destroy AIG \ts{20}{9}{9:18}.

\paragraph{What triggered payment?} Without margin, the premium flows and a large payment comes only at actual default; for AAA tranches, with many tranches beneath,
defaults rarely reached them. It was not default but the fluctuation in CDS values that caused the problem. That is why many underestimated the crisis: ``there
just aren't a lot of subprime mortgages'', ignoring leverage, liquidity and the embedding in world money and capital markets. Those who took the low tranches and
knew they bore risk mostly survived; those who thought they took no risk failed \ts{20}{9}{12:24}.

% ------------------------------------------------------------------
\section{Abacus: the synthetic CDO as collateral prepayment}\label{sec:cds-abacus}

John Paulson wanted to short subprime, expecting the housing bubble to end. He bought CDS on CDO tranches of mortgage-backed securities, insurance without owning the
securities \ts{20}{10}{0:02}. Goldman Sachs set up an entity, Abacus, which took the other side of the CDS, credit exposure as if owning
RMBS, ``just a side bet'' \ts{20}{10}{1:04}. Abacus sold the German bank IKB a CDO tranche paying a spread over the risk-free rate; IKB
borrowed to buy it, and Abacus invested the proceeds in Treasury bills \ts{20}{10}{2:08}.

\begin{nfigure}
\begin{taccts}
\tacct[1.9cm]{Paulson}{CDS on CDO tranches (long protection)}{}\quad
\tacct[1.9cm]{Abacus (synthetic CDO)}{Treasury bills}{CDS sold to Paulson\\CDO tranche sold to IKB}\quad
\tacct[1.9cm]{IKB}{CDO tranche}{Borrowing}
\end{taccts}
\caption{The synthetic CDO: T-bills plus a short CDS replicate a long position in mortgage-backed securities \ts{20}{10}{2:28}.}\label{fig:cds-abacus}
\end{nfigure}

\begin{definition}[New concepts: synthetic CDO]
A \emph{collateralized debt obligation} (CDO) issues tranches backed by a pool of debt; a \emph{synthetic} CDO owns no such debt but replicates the exposure with
risk-free securities plus CDS sold on the reference assets \ts{20}{10}{2:49}.
\end{definition}

When the tranches fell, the CDS rose and the T-bills went to Paulson as margin. From a liquidity perspective that is the genius of it: IKB had paid 100\% upfront,
so there was no problem getting paid, no bill to send to IKB that it could dispute, no need to liquidate RMBS exposure
\ts{20}{10}{3:30}. There were issues about Paulson's role in choosing the tranches most
likely to default and about Goldman's representations, and a fine was paid; but no issue of payment. IKB was bailed out by the German government; AIG's counterparties
included Soci\'et\'e G\'en\'erale: the international aspects again \ts{20}{10}{4:55}.

\begin{history}[Abacus 2007-AC1]
In April 2010 the SEC charged Goldman Sachs with misleading investors in the synthetic CDO Abacus 2007-AC1, whose reference portfolio Paulson \& Co.\ had helped select
while betting against it; Goldman settled in July 2010 for \$550 million. IKB, an early casualty of the crisis in 2007, was rescued with the support of the state-owned
KfW.\par
\emph{References:} SEC press releases 2010-59 (16 April 2010) and 2010-123 (15 July 2010); G.~Zuckerman, \emph{The Greatest Trade Ever}, 2009 (cited in Mehrling's notes, which also
recall, from G.~Tett, \emph{Fool's Gold}, 2009, that the first CDS were built in this synthetic way for J.~P.~Morgan to hedge tail risk on corporate loans) \mn{20}{6}.
\end{history}

% ------------------------------------------------------------------
\flushside
\section*{Key formulas and ideas}
\begin{keyformulas}
\begin{tabularx}{\linewidth}{@{}l L@{}}
Black (1970) & risky bond $+$ IRS $+$ CDS $=$ risk-free bond\\
CDS & notional long Treasury $+$ short corporate bond of the same maturity; buyer long the swap, short credit risk\\
Bond price & $P_0=\sum C_t/(1+r+\delta)^t+F/(1+r+\delta)^T$; moves with spreads and ratings\\
Pricing & bond $+$ IRS $=$ LIBOR $+s$; CDS premium $u$; no-arbitrage $u=s$ if no liquidity premium\\
Negative basis & $u<s$: buy bond, buy protection, fund in money market; liquidity risk in the funding\\
AIG & margin calls on in-the-money CDS: ``liquidity kills you quick''\\
Abacus & synthetic CDO $=$ T-bills $+$ short CDS: collateral prepaid\\
\end{tabularx}
\end{keyformulas}

\section*{Exercises}

\begin{exercise}[Black's decomposition]
Write the balance sheet of an investor who holds a ten-year BBB bond, an interest rate swap and a CDS on it, and show that its net exposure is a Treasury bill. Who
bears the interest rate risk and who the default risk?
\end{exercise}

\begin{exercise}[CDS cash flows]
A protection buyer on \$10 million face pays $u=400$ bp a year. The bond defaults at the end of year 3 and then trades at 35. Write the cash flows each year. What is
the buyer's position worth just before default if spreads have widened to 900 bp?
\end{exercise}

\begin{exercise}[Negative basis]
A AAA tranche yields LIBOR $+$ 60 bp, protection costs 20 bp, and repo funding costs LIBOR $+$ 10 bp. Compute the annual carry on \$1 billion. The tranche's price falls
from 100 to 96 and the repo haircut rises from 2\% to 10\%: how much cash must be found, and from where?
\end{exercise}

\begin{exercise}[Margin or no margin]
Compare an AIG-style CDS seller that posts margin when downgraded with a monoline that posts none, if the marked-to-market value of the CDS rises by \$5 billion but no
default occurs. Which one fails, and why?
\end{exercise}
